%======================================================================
\section{Chern--Simons as symplectic action in the pillowcase}\label{sec:action}
%======================================================================
Normalise $\cs\colon\mathcal B(Z)\to\R/\Z$ so that the energy of an instanton on $\R\times Z$ from $\alpha$ to
$\beta$ is $\cs(\alpha)-\cs(\beta)$ (lifted to $\R$), and so that the trivial connection has $\cs=0$.

\subsection{Forms and primitives}
On $\R^2\ni(a,b)$ let $\omega=2\,da\wedge db$ and $\lambda=2b\,da$.  The factor $2$ makes $\omega$ descend to $P$ with
total area $1$, matching the $\Z$-periodicity of $\cs$.  For an oriented curve $C\subset\R^2$ with a base point
$c_0$, set $f_C(x):=-\int_{c_0}^{x}\lambda$ along $C$.

\begin{lemma}[change of basis]\label{lem:basis}
For $\begin{psmallmatrix}p&q\\r&s\end{psmallmatrix}\in SL(2,\Z)$ and $a'=pa+qb$, $b'=ra+sb$:
$2b'\,da'-2b\,da=dg$ with $g=2qr\,ab+pr\,a^2+qs\,b^2$.  On $L_{p/q}=\{a'=0\}$ parametrised by $(a,b)=s(q,-p)$ one
has $g=pq\,s^2$.
\end{lemma}
\begin{proof}
Expand $b'\,da'=(ra+sb)(p\,da+q\,db)$ and use $ps-qr=1$; the terms regroup as $b\,da+d(qr\,ab+\tfrac12pr\,a^2+\tfrac12qs\,b^2)$.
On the line, substitute.
\end{proof}

\subsection{The Kirk--Klassen formula}
Kirk--Klassen \cite{KK90}: if $\rho_t$ is a path in $R(X)$ with boundary angles $(a(t),b(t))$ lifted continuously to
$\R^2$, and $\rho_0,\rho_1$ extend over the same Dehn filling $Z$, then, in a basis of $H_1(T^2)$ in which the filling
curve is the first basis element,
\[
\cs(\rho_1)-\cs(\rho_0)\equiv-\int_0^1 2b\,\dot a\,dt\pmod{\Z}.
\]
Combining with Lemma~\ref{lem:basis}:

\begin{proposition}[Chern--Simons is the action]\label{prop:action}
Fix $(r,s)$ with $ps-qr=1$ and put $Q(a,b):=pr\,a^2+2qr\,ab+qs\,b^2$ (the $g$ of Lemma~\ref{lem:basis}).  Let
$x,y\in r^{-1}(L_{p/q})$ lie on one component $\gamma$ of $R^{\rm irr}(X)$ and lift a path in $\gamma$ from $x$ to
$y$ to $\tilde\gamma\subset\R^2$ with endpoints $\tilde x,\tilde y$.  Then
\[
\cs(y)-\cs(x)\equiv\big[f_{\tilde\gamma}(\tilde y)-Q(\tilde y)\big]-\big[f_{\tilde\gamma}(\tilde x)-Q(\tilde x)\big]\pmod{\Z}.
\]
Changing $(r,s)$ changes $Q$ by $(pa+qb)^2\in\Z$ at the lifted points, so the right side is well defined.
On $L_{p/q}$ itself, $Q=pq\,s^2$ in the parametrisation $(a,b)=s(q,-p)$; but the lifted endpoints lie on
translates $pa+qb=k$, and $Q$ must be evaluated there, not on the projection to $L_{p/q}$.  For $L_0^w$
take $(p,q)=(0,1)$: $Q=2r\,ab+s\,b^2$ with $s=1$, defined modulo $b^2=\tfrac14$ on $b\in\tfrac12+\Z$.
\end{proposition}
\begin{proof}
Kirk--Klassen in the adapted basis $(a',b')=(pa+qb,\,ra+sb)$, in which the filling curve is the first basis
element and the formula is $-\int2b'\,da'$; rewrite by Lemma~\ref{lem:basis}.  This is the form
given by Andersen--Hansen \cite[eq.~(5.67)]{AH}, who correct a sign in \cite{KK90}.
\end{proof}
The right-hand side is the symplectic action $A(x):=f_{\tilde\gamma}(\tilde x)-Q(\tilde x)$ of the intersection point.
\emph{Numerically verified on all $26$ flat connections of the six homology-sphere surgeries $\pm1,\pm1/2,\pm1/3$ on the
trefoil, and on $\pm1$ surgeries on the figure-eight} (\S\ref{sec:tests}).

\begin{remark}
For $p/q=1/n$ the formula reads $\cs(y)-\cs(x)\equiv-\int 2b\,da+n\,(b(y)^2-b(x)^2)$, which is the form used by
Nozaki--Sato--Taniguchi \cite[Thm.~7.1]{NST} (with their sign convention for the integral).  An integer-valued
lift of the Kirk--Klassen formula, which is what the filtered complex needs, is Boden--Herald--Kirk--Klassen
\cite[Thm.~5.5]{BHKK}: $\cs(A_1)=-c+2\int n\,\dot m\,dt$ with $c$ a linking number of an explicit loop.
\end{remark}

So \emph{the Chern--Simons filtration on the generators of $\CI(Y_{p/q})$ coming from one component of
$R^{\rm irr}(X)$ is the action filtration of the pillowcase Floer complex}.

\begin{lemma}[periods]\label{lem:period}
$f_\gamma$ does not in general descend to $\gamma$.  By \cite{KK93}, $\cs_X$ is a section over $R(X)$ of the circle bundle
$E(\partial X)=\big(\R^2\times U(1)\big)/\!\approx$ with $(a,b;z)\approx(a+1,b;z e^{-8\pi ib})\approx(a,b+1;z e^{8\pi ia})\approx(-a,-b;z)$,
and $e^{2\pi i f_\gamma}$ is the pull-back of that section to the lift.  Hence if one circuit of $\gamma$ translates the lift by
$(m,k)\in\Z^2$, the period $\oint_\gamma 2b\,da$ is an integer only after adding the transition terms of $E(\partial X)$.
\end{lemma}
For the figure-eight knot the lift translates by $(0,\mp2)$ and $\oint2b\,da=\tfrac13$; with the transition term the
corrected period is $-1$ (\S\ref{sec:tests}).  The first draft of this lemma claimed $\oint2b\,da\in\tfrac12\Z$; that is false.

\subsection{Absolute constants: the first correction terms}
Proposition~\ref{prop:action} determines $\cs$ on the generators coming from one component $\gamma_i$ only up to a
constant $c_i\in\R/\Z$: there is no path in $R(X)$ from the trivial connection to $\gamma_i$, precisely because of
\ref{H:alex}.  These constants are genuine knot invariants (they are the Chern--Simons invariants of flat connections
on fillings), and they are computable:
\begin{enumerate}[label=(\roman*)]
\item through the $SL(2,\C)$ character variety, where the Kirk--Klassen formula holds with complex $(a,b)$
\cite{KK93} and the irreducible curve containing $\gamma_i$ reaches the abelian locus at roots of $\Delta_K$ off the
unit circle; the imaginary part of the integral must vanish at the end, which is a check;
\item or from one known filling (e.g.\ a Brieskorn sphere), which calibrates $c_i$ for every slope at once.
\end{enumerate}
Thus the Chern--Simons data of $K$ is: the curve $\Gam$, the primitives $f_{\gamma_i}$, and the constants $c_i$.
